\begin{table*}[!b]
    \centering
    \small
    \begin{tabular}{@{}cccccc@{}}
        \toprule
        $P$ & $S$ & $O$ & \textbf{tokens} & $S \mid P$ & $P \mid S$ \\
        \midrule
        $8$ & $8$ & $0.000$ & $256$ & yes & yes \\
        \midrule
        $16$ & $8$ & $0.500$ & $255$ & yes & \textbf{no} \\
        $16$ & $12$ & $0.250$ & $170$ & \textbf{no} & \textbf{no} \\
        $16$ & $16$ & $0.000$ & $128$ & yes & yes \\
        \midrule
        $24$ & $24$ & $0.000$ & $85$ & yes & yes \\
        \bottomrule
    \end{tabular}
    \caption{The five geometries analysed here, carried over from Deliverable~1. The token count is the number of patch tokens at the $2048$-sample training context. The last two columns are divisibility conditions, and each states that one branch is contained in the other: $S \mid P$ means the patch holds a whole number of strides, so every stride lock is also a patch null, and $P \mid S$ the converse, that every patch null is also a stride lock. Both hold only at $P=S$, where the two branches are the same set of frequencies.}
    \label{tab:patchStride}
\end{table*}
